\section{A constructive real-coordinate schema}\label{sec:real}

The rational theorem already gives a literal iPRA interpretation of the
geometry. This extension explains precisely how the same proofs apply
to constructively given real coordinates, without adding a second-order
domain of arbitrary real numbers.

\subsection{Fixed primitive-recursive names and arithmetic predicates}
Fix, externally, finitely many primitive-recursive rational sequences
$a_i(n)$ with positive denominator functions. Positivity is either built
into the coding, as a denominator of the form $d_i(n)+1$, or included as
an explicit arithmetic premise. Their definitions may have numerical
parameters. Include as hypotheses, or prove in iPRA,
the regularity conditions
\begin{equation}\label{eq:regular}
 |a_i(m)-a_i(n)|\le2^{-m}+2^{-n}\qquad(m,n\in\N).
\end{equation}
Rational inequalities are interpreted by integer cross multiplication of
positive denominators, so these are first-order arithmetic formulas.
The instance at $m=0$ gives $|a_i(n)|\le|a_i(0)|+2$. For example,
$2$ plus the sum of the absolute initial numerators is an integer
bound $B\ge1$ for all the coordinates.

For each fixed integer polynomial $P$ in this finite tuple, write
$P_n=P(a_1(n),\ldots,a_d(n))$ and define
\begin{align}
 \Zero(P)&:\Longleftrightarrow
 \forall k\,\exists N\,\forall n\ge N\ (|P_n|\le2^{-k}),\label{eq:zero-real}\\
 \Apart(P)&:\Longleftrightarrow
 \exists r\,\exists N\,\forall n\ge N\ (|P_n|\ge2^{-r}).\label{eq:apart-real}
\end{align}
In the usual real interpretation, these mean that the polynomial value
is zero, or is positively apart from zero. More fundamentally here, they
are explicit arithmetic formulas. The statement is a schema for each
fixed finite list of PR definitions. It does not quantify over arbitrary
function codes or posit a total universal primitive-recursive evaluator.

\subsection{The estimates replacing exact cancellation}
If $P=\sum_\alpha c_\alpha x^\alpha$, set
\[
 M_P(B)=\sum_\alpha |c_\alpha|B^{|\alpha|}.
\]
Then $|P_n|\le M_P(B)$, term by term. All polynomials in use are fixed
and finite. iPRA proves the following rules from the regularity premises:
\begin{align}
 \Zero(P)\land\Zero(Q)&\Longrightarrow\Zero(P+Q),\label{eq:z-add}\\
 \Zero(P)&\Longrightarrow\Zero(SP),\label{eq:z-mul}\\
 \Apart(C)\land\Zero(CP)&\Longrightarrow\Zero(P).\label{eq:z-cancel}
\end{align}
For \eqref{eq:z-add}, given $k$, use both zero hypotheses at precision
$k+1$ and take the maximum of their supplied thresholds. The triangle
inequality gives $2^{-(k+1)}+2^{-(k+1)}=2^{-k}$.
For \eqref{eq:z-mul}, choose an explicit $t$ with $M_S(B)\le2^t$,
for example $t=M_S(B)+1$, and use the zero hypothesis at precision
$k+t$. The bound on $S$ gives the result.
For \eqref{eq:z-cancel}, open the witnesses $r,N_0$ of apartness.
Use $\Zero(CP)$ at precision $k+r$ and take the maximum of $N_0$ and
the resulting threshold. For every subsequent $n$,
\[
 2^{-r}|P_n|\le|C_nP_n|\le2^{-(k+r)},
\]
hence $|P_n|\le2^{-k}$. This is rational arithmetic with a positive dyadic
factor. Negation preserves $\Zero$, and fixed finite sums are obtained
by finitely many applications of \eqref{eq:z-add}. Apartness of two
factors gives apartness of their product by adding the exponents and
taking the maximum of the thresholds.

Each witness comes from a specified existential premise or a
primitive-recursive operation on finitely many witnesses. No sequence
of witnesses is chosen, and no induction is performed on $\Zero$ or
$\Apart$. Quantifier-free induction suffices for the elementary power
and rational-arithmetic bounds. These are consequently conditional
first-order iPRA derivations even though their hypotheses are quantified.

\subsection{Transferring the geometric proofs}
All polynomial identities hold exactly at every approximation index.
For Desargues, replace incidence through a center by
$\Zero(o\cdot j_i)$ and require $\Apart(o_j)$ for a designated center
coordinate. Equation~\eqref{eq:adj3}, finite sums, bounded multipliers,
and \eqref{eq:z-cancel} give $\Zero(\br{j_1}{j_2}{j_3})$.
The Desargues identity gives $\Zero(\br X Y Z)$.
To produce an actual nonzero line, require apartness of a designated
coordinate of $X\times Y$ and choose $h=X\times Y$. Its first two
incidences are polynomial identities; its third is the zero determinant.

For the converse, the supplied axis satisfies $\Zero(h\cdot X)$,
$\Zero(h\cdot Y)$, and $\Zero(h\cdot Z)$. Require apartness of both
triangle determinants and of a designated coordinate of this axis. The same argument
gives the zero joining-line determinant. Apartness of a designated
coordinate of $j_1\times j_2$ produces the nonzero center
$o=j_1\times j_2$.

For Pascal, assume $\Zero(Q(P_i))$ for all six vertices and
$\Apart(q_j)$ for a designated coefficient. The cofactor identity
\eqref{eq:conic-certificate} gives $\Zero(q_j\det M)$ by finite sums
and bounded multipliers. Cancellation yields $\Zero(\det M)$ and
therefore $\Zero(\br X Y Z)$. A designated apart nonzero coordinate
of $X\times Y$ again gives the Pascal line $h=X\times Y$.
The side and meet vectors have supplied apartness witnesses whenever
they are to denote projective objects.

The rational branch algorithm \eqref{eq:axis} is deliberately not used
on arbitrary real names: exact real equality is not being decided.
For this schema, apartness of a selected pair supplies the pivot.
These are explicit quantitative nondegeneracy assumptions for the
constructive real version; they are stronger than merely negating
equality to zero.

\subsection{Producing regular output names}
A raw polynomial sequence $P_n$ need not satisfy the unit modulus
\eqref{eq:regular}. The repair is explicit. Put
\[
 L_P(B)=\sum_{|\alpha|>0}|c_\alpha|\,|\alpha|B^{|\alpha|-1}.
\]
The telescoping product identity applied to each monomial proves
\[
 |P_m-P_n|\le L_P(B)(2^{-m}+2^{-n}).
\]
Choose $s$ with $L_P(B)\le2^s$, for example $s=L_P(B)+1$.
Then $\widetilde P(n)=P_{n+s}$ is regular in the sense of
\eqref{eq:regular}. For a vector use a common shift for its components.

The same arithmetic formulas $\Zero$ and $\Apart$ also apply to the
derived PR rational sequences below, including shifted differences;
equivalently, enlarge the fixed finite family of names to include their
specified shifts. To see that resampling does not change any incidence, first note
that $P_n-P_{n+s}$ tends to zero with the displayed explicit bound.
For several output components, their respective shifts give
coordinate differences tending to zero; telescoping products and
bounded multipliers then show that every fixed polynomial in the
original outputs and the corresponding polynomial in the resampled
outputs differ by a $\Zero$ sequence. Thus all proven zero equations
transfer. Apartness transfers as well: if $|P_n|\ge2^{-r}$ eventually,
use $|P_n-\widetilde P(n)|\le2^{-(r+1)}$ eventually to obtain
$|\widetilde P(n)|\ge2^{-(r+1)}$. Thresholds are combined by maxima.
Consequently the constructed line or center can be represented by
regular primitive-recursive real names without changing the conclusion.

This extension is a completed conditional proof, not a claim that every
classical real configuration has a primitive-recursive presentation or
that iPRA contains an unrestricted real-number sort.

